\documentclass[11pt,a4paper]{article}
\usepackage[T1]{fontenc}
\usepackage[utf8]{inputenc}
\usepackage{lmodern}
\usepackage[margin=25mm]{geometry}
\usepackage{amsmath,amssymb,amsthm}
\usepackage{array,booktabs,longtable,tabularx}
\usepackage{enumitem}
\usepackage[hidelinks]{hyperref}
\usepackage{microtype}
\setlength{\parindent}{0pt}
\setlength{\parskip}{5pt plus 1pt}
\setlist[itemize]{leftmargin=1.5em,itemsep=2pt,topsep=3pt}
\newcommand{\C}{\mathcal C}
\newcommand{\K}{\mathcal K}
\newcommand{\Lcat}{\mathcal L}
\newcommand{\Rel}{\mathbf{Rel}}
\newcommand{\Corel}{\mathbf{Corel}}
\newcommand{\LinRel}{\mathbf{LinRel}}
\newcommand{\FinSet}{\mathbf{FinSet}}
\newcommand{\Cat}{\mathbf{Cat}}
\newcommand{\op}{\mathrm{op}}
\newcommand{\id}{\mathrm{id}}
\newcommand{\im}{\operatorname{im}}
\newcommand{\dom}{\operatorname{dom}}
\newcommand{\cldf}{\operatorname{cl}_{\mathrm{df}}}
\title{Mixed context and content normal forms\\[3pt]\large A bounded reduction test}
\author{}
\date{1 November 2026}
\begin{document}
\maketitle
\vspace{-1.5em}

\textbf{Conclusion.} The finite-set candidate fails before the usual coherence
question: its crossing gives different answers on two representatives identified
by a pure middle map. It therefore does not descend to the balanced tensor of
the existing one-sided calculi. Its individual normal forms and squares have an
exact corelational description, but the resulting normalization is not compatible
with repeated composition. Ordinary relations preserve the partial-map control
while losing content distinctions; corelations retain those distinctions while
changing partial-map composition. The linear control is the established calculus
of linear relations. Generic composition from a supplied distributive law is
already covered; the residual is an independently obtained balanced crossing with
the stated preservation requirements, plus the unproved full weak retained-cell
comparison.

\textbf{Scope.} This is the focused reduction test authorized after the scoping
audit, not execution of the whole provisional research contract. It stops at
precise covered cases and a residual obligation. No new theorem package is
adopted, no novelty is promised, and no impossibility claim is made about every
mixed contextad construction. The broader programme remains diverse theoretical
and applied generalizations; robotics is only one possible application.

\section{The typed target and the required comparison}

Retain the prior indexed inputs, with their supplied one-sided structures
\cite{CM},
\(\K:\C^{\op}\to\Cat\), \(\Lcat:\C\to\Cat\), with context and content
extensions \(e_AP\), \(d_BQ\). A mixed arrow is
\[
 (P,Q,f):A\rightsquigarrow B,
 \qquad P\in\K_A,\ Q\in\Lcat_B,\quad f:e_AP\to d_BQ.
\]
At fixed endpoints, a comparison consists of
\(\alpha:P\to P'\), \(\beta:Q\to Q'\) with
\[
 d_B(\beta)f=f'e_A(\alpha).
 \tag{1}
\]
Squares over base maps \(u:A\to A'\), \(v:B\to B'\) use
\(\alpha:P\to u^*P'\), \(\beta:v_!Q\to Q'\) and the supplied transport maps.
These directions matter: the hom-category is a comma category, not merely a set
of isomorphism classes. Tight maps remain the original base maps. Their pure
horizontal representative is
\((I_A,J_B,\eta_Bu\varepsilon_A)\).

The missing general interaction is still typed at the correct base. For
\(Q\in\Lcat_B\), \(R\in\K_B\), it needs
\[
 \widehat R\in\K_{d_BQ},\quad
 \widehat Q\in\Lcat_{e_BR},\quad
 \chi:e_{d_BQ}\widehat R\to d_{e_BR}\widehat Q.
 \tag{2}
\]
Writing \(f^*R\) directly is ill-typed. Once (2) is supplied, the composite
uses \(P\otimes f^*\widehat R\), \(S\diamond g_!\widehat Q\) and the
previously displayed extension--transport--crossing formula. Its typing does
not establish its coherence.

A reduction must distinguish preservation of individual arrows, fixed-endpoint
cells, squares over base maps, horizontal composition, both one-sided controls,
and dependence. Equality after discarding witnesses is not the same assertion
as an equivalence of these retained categories. The tests below do not replace
arbitrary dependent families by a finite-set example without saying so.

\section{The finite-set normal forms are exactly presented difunctional relations}

For the display/codisplay test, take monomorphisms and epimorphisms in
\(\FinSet\). Up to isomorphism of quotient labels, a normal form is
\[
 D\subseteq A,\qquad q:B\twoheadrightarrow Y,\qquad f:D\to Y.
 \tag{3}
\]
Its relation is
\[
 R_f(a,b)\quad\Longleftrightarrow\quad a\in D\ \text{and}\ f(a)=q(b).
 \tag{4}
\]
Every nonempty row is a complete quotient block. Thus nonempty rows are equal
or disjoint, precisely the difunctionality condition
\(aRb,\ a'Rb,\ a'Rb'\Rightarrow aRb'\).
Conversely, any finite difunctional relation determines its domain and its
nonempty row blocks; partitioning the unused target subset arbitrarily completes
(3). Therefore the normal form is a difunctional relation \emph{plus a partition
of the unused target}. Those extra blocks are not just quotient labels.

For a following form \(E\subseteq B\), \(r:C\twoheadrightarrow Z\),
\(g:E\to Z\), the proposed recipe takes
\[
 W=q(E),\quad D'=f^{-1}(W),\quad W+_EZ,
 \tag{5}
\]
using \(q|_E\) and \(g\) for the pushout. Define the expanded relation
\[
 M(y,c)\quad\Longleftrightarrow\quad
 \exists b\in E:\ q(b)=y,\quad g(b)=r(c).
\]
It includes active quotient nodes which no first input selects. Let
\(\cldf\) complete every nonempty connected component of a bipartite relation
to a complete bipartite relation. This is the least difunctional superrelation.
The exact relation of (5) is
\[
 R_{\mathrm{norm}}(a,c)
 \quad\Longleftrightarrow\quad
 a\in D\ \text{and}\ \cldf(M)(f(a),c).
 \tag{6}
\]
Indeed, compressing each complete \(r\)-block in the graph of \(M\) gives
exactly the pushout graph \(W+Z\), with edges \(q(b)\sim g(b)\).
The connected-component test is therefore the output equality test in (5).

If every active node is selected, \(W\subseteq f(D)\), replacing quotient
nodes by their nonempty sets of selecting inputs preserves connectivity. In this
\emph{active-tight} case only,
\[
 R_{\mathrm{norm}}=\cldf(R_gR_f).
 \tag{7}
\]
In general \(R_gR_f\subseteq\cldf(R_gR_f)\subseteq R_{\mathrm{norm}}\),
and the second inclusion can be strict. Formula (6), not the unrestricted
version of (7), is the reduction of the recipe.

\subsection{Unused content already obstructs binary descent}

Let \(A=\{a\}\), \(B=\{0,1,2\}\), \(D=A\), and select the block \(\{0\}\).
Compare partitions
\[
 \{0\}\mid\{1\}\mid\{2\},\qquad \{0\}\mid\{1,2\}.
\]
Both denote \(\{(a,0)\}\), and the first has a noninvertible quotient-coarsening
cell to the second. Follow by the pure function
\(0\mapsto c_0,\ 1\mapsto c_0,\ 2\mapsto c_1\).
Recipe (5) gives the selected row \(\{c_0\}\) in the first case and
\(\{c_0,c_1\}\) in the second: the unused block joins the two outputs.
Ordinary relational composition gives only \(\{c_0\}\) in both cases.
Thus forgetting unused blocks is not a congruence for this binary recipe,
even under pure postcomposition. This also witnesses failure to preserve the
contentful, quotient-pushout control, since both contexts are total.

\subsection{The earlier associativity failure is premature closure}

The first arrow has rows \(\{b_0,b_1\}\), \(\{b_2,b_3\}\) at
\(a_0,a_1\). The next pure function sends these four elements respectively to
\(c_0,c_1,c_1,c_2\). Its ordinary relational composite has rows
\[
 \{c_0,c_1\},\qquad \{c_1,c_2\}.
\]
They overlap without being equal, so they cannot be normal-form rows.
Normalization closes both to \(\{c_0,c_1,c_2\}\). A final partial map defined
only at \(c_2\), with value \(*\), consequently accepts both inputs. In the
opposite grouping the last two maps first restrict to \(\{b_3\}\), so only
\(a_1\) survives. The domains differ even over \(A\); changing representatives
or adding invertible associators cannot repair this recipe. Ordinary \(\Rel\)
retains the non-difunctional intermediate and gives \(\{(a_1,*)\}\) in both
groupings.

\section{Cells and the stronger corelational reduction}

At fixed finite endpoints, presentations form a preorder: a cell exists exactly
when the domain grows, the target partition coarsens, and selected blocks map
to the corresponding selected blocks. Let \(E\) forget to difunctional
relations ordered by inclusion, and \(F\) use singleton unused target blocks.
Then
\[
 F\dashv E,\qquad EF=\id,
 \qquad F(R)\leq P\ \Longleftrightarrow\ R\subseteq E(P).
 \tag{8}
\]
This is a local hom-category statement. Difunctional relations are not closed
under ordinary composition, and (8) does not supply a horizontal functor.

There is a stronger faithful description. A finite-set corelation is a partition
of \(A+B\). Embed (3) by putting each target block together with all source
inputs selecting it, and make each undefined source input a singleton block.
The image consists exactly of partitions whose source-only blocks are
singletons. This embedding \(I\) preserves and reflects all presentation cells,
including squares over arbitrary endpoint maps; it retains unused target blocks.
It is a description of individual arrows and squares, not yet composition.

Let \(N\) split every source-only block of any corelation into singletons,
leaving all other blocks unchanged. For the hom-preorders,
\[
 I\dashv N,\qquad NI=\id.
 \tag{9}
\]
A partition in the image refines a corelation exactly when it refines this
splitting: a block meeting the target cannot land inside a source-only block,
and a source-only image block is already a singleton.
The exact presentation-level formula is
\[
 I(P_2\star P_1)=IN\bigl(I(P_2)\circ_{\Corel}I(P_1)\bigr),
 \tag{10}
\]
where \(N\) returns the normalized presentation. Composition of
corelations joins equivalence classes across the middle boundary; splitting its
source-only classes is precisely the information loss made by (5).
The coreflection is not compatible with repeated horizontal composition.
Equation (10) identifies the normalization inside known machinery, rather than
proving it associative.

\textbf{Why keeping all corelations is not a conservative repair.}
A pure map sending two distinct inputs \(b_1,b_2\) to the same \(c\), followed
by a partial map undefined at \(c\), is nowhere defined on those inputs.
Its normal partial-map presentation has two separate undefined-source
singletons. Corelational composition instead retains the source-only block
\(\{b_1,b_2\}\). Hence \(\Corel\) preserves the total-context contentful
cospan control but not the partial-map contextful control. This failure already
uses a pure map followed by a partial map; retaining more witnesses changes
which one-sided composition is recovered.

\section{The linear positive control is already covered}

For finite-dimensional vector spaces over a field, use a subspace \(D\leq A\), a surjection
\(q:B\twoheadrightarrow Y\), and a linear \(f:D\to Y\). Their relation is
the linear subspace
\[
 R=\{(a,b):a\in D,\ f(a)=q(b)\}\leq A\oplus B.
\]
Conversely a linear relation determines
\(D=\operatorname{pr}_A(R)\), \(K=\{b:(0,b)\in R\}\), the quotient
\(B\to B/K\), and \(f(a)=b+K\) whenever \((a,b)\in R\).
There is no independent unused-block partition. Presentation cells correspond
to relation inclusion, with the analogous endpoint-square condition.

In the linear pushout \(W+_EZ\), equality between a \(W\)-element and a
\(Z\)-element is witnessed by one \(e\in E\): sums of gluing witnesses are
again one linear witness. Therefore (5) agrees with ordinary linear-relational
composition, not a premature closure. This is the existing linear-relation
control; Fong--Zanasi's abelian relation/corelation comparison and linear
specialization are the relevant stronger baseline \cite{FZ}. It establishes a
positive specialization, not a novel mixed calculus or a comparison for arbitrary
dependent families.

\section{The general existing-theorem bridge}

There are two exact reductions, at different preservation levels. First, let
\(e:K\to\C\), \(d:L\to\C\) be functors between small categories. Ordinary
postcomposition and precomposition give strict functors in the intermediate
object. Their ordinary \(\Cat\)-valued coend satisfies
\[
 \int^{Z\in\C}(e\downarrow Z)\times(Z\downarrow d)
 \ \cong\ (e\downarrow d).
 \tag{11}
\]
The map composes \(eP\xrightarrow{x}Z\xrightarrow{y}dQ\).
The inverse represents \(f:eP\to dQ\) at \(Z=eP\), with first leg \(\id\).
For a cell \((\alpha,\beta)\), first apply \(\beta\) there, balance along
\(e\alpha\), then apply \(\alpha\) in the first factor. Product-category
interchange proves functoriality; balancing verifies both inverse identities on
morphisms as well as objects. With \(e=e_A,d=d_B\), (11) retains exactly the
mixed cells (1), including noninvertible ones. It is local: no horizontal or
pseudoambient coherence theorem follows just from this calculation.

Second, pass the supplied one-sided horizontal hom-categories to
\emph{isomorphism classes of arrows}, not connected components. They give ordinary
categories \(M,N\) and identity-on-objects pure functors
\(J_M:\C\to M\), \(J_N:\C\to N\). Such a functor is equivalently a
promonad in ordinary \(\mathbf{Prof}\): its profunctor is the arrow set,
its actions compose with pure maps, and its multiplication is one-sided
composition. Rom\'an's Theorem 3.9, with explicit proof in Theorem B.2, is a
precise existing baseline \cite{Roman}.
Using left-to-right execution order,
\[
 (M\diamond N)(A,B)=\int^Z M(A,Z)\times N(Z,B)
\]
is bijective to mixed triples modulo grade isomorphisms, by the same splitting
at \(e_AP\). A distributive law
\(\tau:N\diamond M\to M\diamond N\) gives the composite promonad by Beck's
standard construction. Its value on the canonical content arrow
\(d_BQ\rightsquigarrow B\) and context arrow
\(B\rightsquigarrow e_BR\), each with identity body, is an isomorphism class of triples of shape (2).
Choosing a representative is not yet a coherent lift retaining cells.
This is an actual ordinary reduction, but it discards noninvertible cells.
The full weak cell-preserving ambient comparison has not been established.

\textbf{The first failed condition is balancing.}
In the three-stage example, write \(f\) for the total-context contentful arrow,
\(g\) for the middle pure map, and \(h\) for the final contextful partial map.
The source tensor \((N\diamond M)(A,D)\) identifies
\[
 [\,f;J_N(g),h\,]=[\,f,J_M(g);h\,],
 \tag{12}
\]
where semicolons denote one-sided composition in execution order. The left
representative first uses the native contentful pushout; the right first uses
the native contextful pullback. The tested crossing gives domain
\(\{a_0,a_1\}\) on the left and \(\{a_1\}\) on the right.
It therefore does not even define the proposed map \(\tau\), including modulo
source-preserving isomorphism. This is sharper than failure of a later Beck
multiplication axiom or associator, and concerns this crossing, not all possible
products.

Street's mixed opwreath construction (Definition 5 and Proposition 7) and
Garner's two-sided Kleisli construction (Definition 14) already provide the
formal fixed-ambient mixed pattern \cite{Street,Garner}. Garner also derives a
nontrivial law on \(\mathrm{Mod}\), so deriving laws is not itself an unexplored
generic problem. Walker's law--lifting equivalences offer alternate presentations,
not an automatic source of the missing interaction \cite{Walker}. None supplies
an unverified dependent crossing merely from the separate one-sided inputs.

\section{Preservation and the residual question}

\small
\begin{longtable}{@{}p{0.20\linewidth}p{0.24\linewidth}p{0.23\linewidth}p{0.23\linewidth}@{}}
\toprule
Comparison & Individual arrows and squares & Horizontal and one-sided controls & Dependence retained \\
\midrule
\endhead
Finite forms to \(\Rel\) & Loses unused target partitions; cells map to inclusions & Not the proposed recipe; partial maps preserved, contentful pushouts not preserved & Result-selected subdomains remain; not arbitrary indexed families \\
Finite forms into \(\Corel\) & Fully faithful on the presentation arrow and square data & Recipe is normalized composition; full \(\Corel\) fails partial-map control & Retains the displayed finite data, not a general dependent bridge \\
Linear forms to \(\LinRel\) & Equivalence up to presentation isomorphism, retaining cells & Recipe and both one-sided controls preserved & Subspace-display specialization only \\
General bridge & Strict local coend retains cells; ordinary Prof reduction discards noninvertible cells & Tested crossing fails pure balancing before Beck axioms & Local dependent grades retained; full weak transport comparison unproved \\
\bottomrule
\end{longtable}
\normalsize

Pure base maps remain identified in each individual-arrow comparison above.
That fact alone does not imply preservation of composites mixing pure maps
with effects, as the two finite counterexamples show.

The bounded test therefore stops with two covered descriptions and one precise
residual. The linear calculus is an established positive case; finite normal
forms are a locally fully faithful subpresentation of corelations with a
nonmultiplicative normalization. Neither observation yields the desired general
dependent mixed theorem. The residual is an existence, obstruction or independent-criterion theorem for
a balanced crossing of specified indexed inputs, conservative on both one-sided
calculi. A positive result must first prove pure-map balancing and endpoint
naturality, then verify the existing unit and multiplication axioms and establish the
cell-preserving comparison and both one-sided recoveries. Merely assuming that
law would be a useful specialization theorem, not a new entry criterion.

No adequate general continuation restriction or universal repair has been
established here. Choosing relations, corelations, a linear-only scope, or a
separately specified continuation interface changes a preservation commitment
and remains a scope decision. The existing contract and its review gate remain
unchanged.

\section{Evidence and verification limits}

The companion ZIP contains the finite-family proofs, the retained-cell and
corelation proof, the exact general bridge analysis, targeted primary-source
references, and small standard-library checkers. The relational checker tests
all 1,645 normal-form pairs on carrier sizes \(2\to3\to2\), including 610
active-tight pairs, and reconstructs all 34 difunctional relations on the
\(2\times3\) carrier. The cell checker verifies the stated bounded adjunction,
embedding and normalization identities. The symbolic arguments, not enumeration,
support the family claims.

This is a targeted reduction, not an exhaustive literature or proof-dependency
certification. Source proof-status caveats in the prior audit remain in force.
The package excludes primary-source PDFs and extracted source text. The TeX is
textually and structurally checked but not compiled or visually rendered, as
permitted.

\begin{thebibliography}{9}
\bibitem{CM} Matteo Capucci and David Jaz Myers.
\emph{Contextads as Wreaths; Kleisli, Para, and Span Constructions as Wreath Products}.
arXiv:2410.21889v1, 2024.
\url{https://arxiv.org/abs/2410.21889}.
\bibitem{Street} Ross Street.
\emph{Wreaths, mixed wreaths and twisted coactions}.
arXiv:1610.07281.
\url{https://arxiv.org/abs/1610.07281}.
\bibitem{Garner} Richard Garner.
\emph{Polycategories via pseudo-distributive laws}.
arXiv:math/0606735v2.
\url{https://arxiv.org/abs/math/0606735v2}.
\bibitem{Walker} Charles Walker.
\emph{Presentations of pseudodistributive laws}.
arXiv:2102.12468v4.
\url{https://arxiv.org/abs/2102.12468v4}.
\bibitem{Roman} Mario Rom\'an.
\emph{Promonads and String Diagrams for Effectful Categories}.
arXiv:2205.07664v4, 1 May 2024, Theorem 3.9 and Theorem B.2.
\url{https://arxiv.org/abs/2205.07664v4}.
\bibitem{FZ} Brendan Fong and Fabio Zanasi.
\emph{Universal Constructions for (Co)Relations: categories, monoidal categories, and props}.
Extended paper, arXiv:1710.03894, especially Corollary 3.4 and Section 4.7.
\url{https://arxiv.org/abs/1710.03894}.
\end{thebibliography}
\end{document}
